\subsection{Limitations and Possible Improvements}\label{sec:limitations}
\paragraph{Limitations of the solution.}
\begin{itemize}
      \item The route is approximate (1.54\% above optimal), and it depends on the
            random seed: the 100-seed results in Table~\ref{tab:experiments} range from
            13.61 to 14.26 km.
  \item Distance is the only cost. Traffic, travel time, opening hours and the time spent at each stop
        are ignored, although they matter for a real one-day tour.
  \item Weights are made symmetric by averaging, so one-way streets are only approximated (this route is
        13.54 km on the raw one-way values).
  \item Road distances come from OpenStreetMap routing. A Google Maps check of
        the final route gave 14.40 km against 13.82 km (+4.2\%), with individual
        legs differing by up to 39\% (Table~\ref{tab:gmaps}). The optimality
        statement of Section~5.2 holds for the OSRM matrix only; we did not
        recompute the optimum with Google distances.
\end{itemize}

\paragraph{Improvements.}
\begin{itemize}
  \item \textbf{Local search after merging.} Applying 2-opt or Or-opt moves to
        the final tour removes locally poor orderings and usually closes most of
        the remaining gap~\cite{applegate2006tsp}. On our route, a single 2-opt
        move that reverses the segment $V_4 \to V_5$ replaces the edges
        $V_3$--$V_4$ (1.31 km) and $V_5$--$V_6$ (0.74 km) by $V_3$--$V_5$
        (0.97 km) and $V_4$--$V_6$ (0.87 km). This saves $2.05 - 1.84 = 0.21$
        km and gives exactly the optimal length of 13.61 km.
  \item \textbf{Multiple restarts.} Run several seeds and keep the shortest tour. Each run takes about
        0.3~ms here, and the best of \RunCount{} runs was \BridgeBest{}~km. This equals the optimum
        only because the separate Held--Karp computation identifies it; restarts alone do not prove
        optimality.
  \item \textbf{Merge rule.} We also tested the best-splice rule, which may
        break any edge of either sub-tour rather than only the edges next to the
        shortest bridge. On this instance it reached the optimum more often (17 vs.\
        6 runs), but its mean length was 13.91 km against 13.88 km
        (Table~\ref{tab:experiments}). A cheaper merge at each step does not
        guarantee a shorter final tour, because every merge is still a local,
        irreversible decision.
  \item \textbf{Richer weights.} Use travel time with traffic from Google Maps, or a directed (asymmetric)
        matrix, if the objective changes from distance to time.
  \item \textbf{Exact methods for small instances.} For about 15 stops or fewer, Held--Karp is fast
        enough to give the true optimum. Divide and Conquer is most useful when $n$ is too large for that.
\end{itemize}
